\chapter{Conclusiones}
\label{chap:conclusions}

\lettrine{E}{l} desarrollo de PadelCenter ha permitido poner en práctica, en un
proyecto real de alcance completo, las competencias adquiridas durante el Grado
en Ingeniería Informática, mención Ingeniería del Software. En este capítulo se
evalúan los resultados obtenidos, se analizan las dificultades encontradas y se
proponen líneas de trabajo futuro.

\section{Objetivos alcanzados}
\label{sec:objetivos_alcanzados}

Todos los objetivos planteados en la sección~\ref{sec:objetivos} han sido
satisfactoriamente alcanzados:

\begin{itemize}
  \item Se ha implementado un sistema completo de autenticación JWT con refresh
    token via cookie \textit{httpOnly}, cobertura de roles y protección de rutas
    tanto en el backend como en el frontend.

  \item El módulo de gestión de centros, pistas y reservas funciona correctamente,
    incluyendo el control de conflictos temporales y la diferenciación de
    visibilidad según el rol del usuario.

  \item El módulo de torneos permite la inscripción de jugadores, la generación
    automática de emparejamientos y el registro de resultados.

  \item La arquitectura hexagonal se ha aplicado de forma rigurosa y se verifica
    automáticamente en cada build mediante ArchUnit.

  \item El enfoque \textit{contract-first} con OpenAPI es efectivo: toda la API
    está documentada en el YAML y tanto los interfaces del backend como los hooks
    del frontend se generan automáticamente a partir de él.

  \item El entorno de desarrollo es completamente reproducible con Docker Compose.
\end{itemize}

\section{Dificultades encontradas}
\label{sec:dificultades}

\paragraph{Integración Orval con App Router de Next.js} La generación de clientes
HTTP con Orval no es trivial cuando se combina con el App Router de Next.js~14,
que distingue entre \textit{Server Components} y \textit{Client Components}.
Los hooks React Query solo pueden usarse en componentes de cliente, lo que
requiere una estrategia de prefetch en el servidor y uso de \texttt{use client}
en los componentes que consumen los hooks.

\paragraph{Control transaccional en tests de integración} Los tests de integración
con Testcontainers y Spring requieren un cuidado especial en el aislamiento de
datos entre tests. El uso de \texttt{@Transactional} en los tests puede enmascarar
problemas reales con las transacciones de la aplicación, por lo que se optó por
truncar las tablas entre tests de forma explícita.

\paragraph{Generación de emparejamientos} El algoritmo de generación de
emparejamientos para torneos con número impar de participantes requirió un
diseño cuidadoso para gestionar los \textit{bye} (participantes con pase
automático a la siguiente ronda) de forma consistente.

\section{Relación con las competencias de la titulación}
\label{sec:competencias}

El desarrollo de PadelCenter ha permitido adquirir y consolidar las siguientes
competencias del GEI, mención Ingeniería del Software:

\begin{description}
  \item[CE1 — Accesibilidad, usabilidad y seguridad:]\leavevmode
    El diseño de la API REST, la implementación de la autenticación JWT y la
    protección de rutas demuestran la aplicación de esta competencia.

  \item[CE2 — Plataformas hardware y software:]\leavevmode
    La evaluación comparativa de frameworks (Spring Boot, Quarkus, Next.js) y
    la justificación de las elecciones tecnológicas refleja esta competencia.

  \item[CE3 — Calidad e innovación tecnológica:]\leavevmode
    La estrategia de pruebas (unitarios, integración, arquitectura) y la
    verificación automática de las reglas de arquitectura son manifestaciones
    de esta competencia.

  \item[CE4 — Servicios y sistemas software:]\leavevmode
    El ciclo completo desde el análisis de requisitos hasta la implementación
    y validación mediante tests demuestra esta competencia.

  \item[CE5 (mención) — Especificación, diseño e implementación:]\leavevmode
    La aplicación de arquitectura hexagonal, el enfoque \textit{contract-first}
    y el uso de herramientas modernas de generación de código son ejemplos
    concretos de esta competencia aplicada.
\end{description}

\section{Líneas de trabajo futuro}
\label{sec:lineas_futuras}

PadelCenter es un sistema funcional y extensible, pero quedan varias líneas de
trabajo que mejorarían significativamente el producto:

\begin{description}
  \item[Autenticación OAuth2/OIDC:] Integración con un proveedor de identidad
    como Keycloak o Auth0, permitiendo el inicio de sesión con Google u otras
    cuentas sociales. La arquitectura hexagonal facilita este cambio: se sustituye
    el adaptador JWT por un adaptador OAuth2 sin tocar el dominio.

  \item[Notificaciones push:] Envío de recordatorios de reserva y notificaciones
    de torneos mediante WebSockets o Server-Sent Events, complementado con
    notificaciones push en dispositivos móviles.

  \item[Aplicación móvil:] Desarrollo de una aplicación React Native que reutilice
    los mismos hooks generados por Orval, ya que la especificación OpenAPI también
    puede generar clientes para React Native.

  \item[Despliegue en cloud:] Contenerización completa con Docker y despliegue
    en un proveedor cloud (AWS, GCP o Azure) mediante Kubernetes o un servicio
    PaaS. La base de datos se migraría a un servicio gestionado (RDS, Cloud SQL).

  \item[Sistema de valoraciones:] Los jugadores podrían valorar las pistas y los
    centros tras una reserva, añadiendo un módulo de reputación al sistema.

  \item[Integración de pagos:] Integración con una pasarela de pago (Stripe) para
    gestionar el cobro de reservas de forma automática, incluyendo reembolsos
    en caso de cancelación.
\end{description}
